% ============================================================
% APPENDIX: Balance-of-Payments Decomposition and Leverage
%           Ratio Algebra
% Cross-referenced from Section C (eq. levr_lom, eq. binding_condition)
% Contains: Steps 1–6 of the formalization
% ============================================================


\section{The Leverage Ratio Embedded in the Balance-of-Payments Constraint}
\label{app:bop_levr}

This appendix develops the full accounting decomposition underlying the
leverage ratio's law of motion (equation~\ref{eq:levr_lom}) and the
constraint-binding condition (equation~\ref{eq:binding_condition})
presented in Section~\ref{sec:C2_bop_dynamics}. The derivation proceeds
in six steps: definition of the leverage ratio, decomposition of the
current and capital accounts, the general balance-of-payments identity,
substitution into the leverage ratio, and a reading of the constraint
under alternative scenarios.

\subsection{The leverage ratio in nominal terms}
\label{app:levr_nominal}

Let all variables be denominated in current domestic currency (CLP) at
time $t$. The exchange rate convention is $e_t$ = CLP per USD, so that
depreciation implies $e_t \uparrow$.

Define:
\begin{itemize}
\item $IR_t$: stock of international reserves, denominated in USD
\item $e_t$: nominal exchange rate (CLP/USD)
\item $M_t$: currency in circulation---the central bank's active domestic
  monetary liability (CLP)
\end{itemize}

The central bank leverage ratio in nominal terms:
%
\begin{equation}
\label{app:eq:levr_def}
\widetilde{\text{LevR}}_t \;\equiv\; \frac{e_t \cdot IR_t}{M_t}
\end{equation}
%
This is a dimensionless ratio: the peso value of the central bank's
world-money holdings relative to its domestic monetary liability issuance.
In growth rates:
%
\begin{equation}
\label{app:eq:levr_growth}
g_{\widetilde{\text{LevR}},t} \;=\; g_{e,t} \;+\; g_{IR,t} \;-\; g_{M,t}
\end{equation}


\subsection{Current account decomposition}
\label{app:current_account}

The current account in domestic currency:
%
\begin{equation}
\label{app:eq:CA}
CA_t = e_t \cdot P_{M,t}^{USD}\!\left(\text{ToT}_t \cdot X_t
       - M_t^{goods}\right)
       - e_t \cdot r_t \cdot D_t^{ext} + T_t
\end{equation}
%
where $\text{ToT}_t \equiv P_{X,t}^{USD} / P_{M,t}^{USD}$ is the net
barter terms of trade. The three components are:

\begin{table}[h]
\centering
\small
\caption{Components of the Current Account}
\begin{tabular}{lll}
\toprule
Term & Content & Character \\
\midrule
$e_t P_{M,t}^{USD}(\text{ToT}_t \cdot X_t - M_t^{goods})$
  & Trade balance at import prices
  & ToT exogenous for small open economy \\
$e_t \cdot r_t \cdot D_t^{ext}$
  & External debt service
  & Largely predetermined \\
$T_t$
  & Net current transfers
  & Exogenous \\
\bottomrule
\end{tabular}
\label{app:tab:CA}
\end{table}

The trade balance is expressed at import prices so that the terms-of-trade
effect enters multiplicatively through $\text{ToT}_t$. A terms-of-trade
improvement ($\text{ToT}_t \uparrow$) raises the trade balance for given
export and import volumes. Debt service is largely predetermined by the
existing stock of external obligations $D_t^{ext}$ and the interest rate
$r_t$ on external debt. Net transfers $T_t$ are exogenous.


\subsection{Capital account decomposition}
\label{app:capital_account}

The capital account in domestic currency:
%
\begin{equation}
\label{app:eq:KA}
KA_t = e_t \cdot \Delta D_t^{ext} + e_t \cdot FDI_t + e_t \cdot \Pi_t
\end{equation}
%
where $\Delta D_t^{ext}$ is net external borrowing, $FDI_t$ is net foreign
direct investment, and $\Pi_t$ denotes net portfolio and short-term capital
flows, with $\Pi_t < 0$ representing capital flight. The capital account
is the channel through which the profit squeeze translates into reserve
drainage: when the domestic profit rate falls and the external constraint
is simultaneously tightening, domestic capital has both the motive (low
profitability) and the means (still-available foreign exchange) to move
resources abroad, compounding the leverage deterioration from the current
account side.


\subsection{The general balance-of-payments identity}
\label{app:bop_identity}

The change in the peso value of international reserves equals the sum of
the current and capital account balances:
%
\begin{equation}
\label{app:eq:bop}
\Delta(e_t \cdot IR_t) = CA_t + KA_t \equiv \text{BoP}_t
\end{equation}
%
Three structural channels operate through this identity:

\begin{enumerate}
\item \textbf{Trade channel:} Terms-of-trade deterioration compresses
  export revenue relative to imports, draining reserves through the
  current account.
\item \textbf{Financial channel:} Debt service outflows or capital flight
  drain reserves independently of trade performance.
\item \textbf{Valuation channel:} Exchange rate depreciation raises the
  CLP value of existing USD reserves---a pure balance sheet effect that
  changes $\widetilde{\text{LevR}}_t$ without altering the underlying
  real position.
\end{enumerate}


\subsection{Substitution into the leverage ratio}
\label{app:levr_substitution}

From the definition $e_t \cdot IR_t = \widetilde{\text{LevR}}_t \cdot
M_t$, substituting into equation~\eqref{app:eq:bop}:
%
\begin{equation}
\label{app:eq:levr_sub}
\Delta\!\left(\widetilde{\text{LevR}}_t \cdot M_t\right) = CA_t + KA_t
\end{equation}
%
Expanding the left-hand side using the product rule and rearranging:
%
\begin{equation}
\label{app:eq:levr_lom}
\boxed{\;
\Delta\widetilde{\text{LevR}}_t \;=\;
\frac{CA_t + KA_t}{M_t}
\;-\; \widetilde{\text{LevR}}_t \cdot g_{M,t}
\;}
\end{equation}
%
This is equation~\eqref{eq:levr_lom} in the main text. The first term is
the net balance-of-payments flow per unit of domestic monetary liability.
The second term---$\widetilde{\text{LevR}}_t \cdot g_{M,t}$---is the
monetary expansion multiplier: the existing leverage ratio scaled by the
growth rate of domestic money. The product structure implies that the same
rate of monetary expansion ($g_{M,t}$) produces a larger absolute drain on
the leverage position when $\widetilde{\text{LevR}}_t$ is high, and a
smaller absolute drain when it is low---but in the latter case the margin
of safety is also smaller, so that even modest monetary expansion can
push the ratio below the binding threshold.


\subsection{Reading the constraint: cases and scenarios}
\label{app:constraint_cases}

The leverage position is deteriorating whenever:
%
\begin{equation}
\label{app:eq:binding}
\Delta\widetilde{\text{LevR}}_t < 0 \;\iff\;
\frac{CA_t + KA_t}{M_t} < \widetilde{\text{LevR}}_t \cdot g_{M,t}
\end{equation}
%
This is equation~\eqref{eq:binding_condition} in the main text. A falling
leverage ratio does not mean the constraint is binding---it means the
constraint is tightening, the distance to the binding threshold is
shrinking. Whether the constraint actually binds depends on the
\textit{level} of $\widetilde{\text{LevR}}_t$.

\paragraph{Case 1: Balance-of-payments deficit}
($CA_t + KA_t < 0$). The constraint tightens trivially. Two forces
compound: reserves are depleted by the BoP deficit (numerator drains),
and monetary expansion erodes the coverage ratio independently
(denominator dilutes).

\paragraph{Case 2: Balance-of-payments surplus}
($CA_t + KA_t > 0$). The constraint tightens whenever the surplus is
insufficient to cover the monetary expansion multiplier:

\begin{table}[h]
\centering
\small
\begin{tabular}{ll}
\toprule
Condition & $\Delta\widetilde{\text{LevR}}_t$ \\
\midrule
Surplus $<$ multiplier threshold & $< 0$ --- tightening despite surplus \\
Surplus $=$ multiplier threshold & $= 0$ --- stable \\
Surplus $>$ multiplier threshold & $> 0$ --- loosening \\
\bottomrule
\end{tabular}
\caption{Leverage dynamics under balance-of-payments surplus.}
\label{app:tab:surplus}
\end{table}

\paragraph{Historical scenarios.}
Three configurations illustrate the constraint's operation across the
regimes that the local projections specification in
Section~\ref{sec:C2_bop_dynamics} conditions on:

\begin{table}[h]
\centering
\small
\begin{tabular}{lccll}
\toprule
Scenario & $e_t \cdot IR_t$ & $M_t$ &
  $\Delta\widetilde{\text{LevR}}_t$ & Mechanism \\
\midrule
A --- Fiscal monetisation
  & Stable & Rising & Falling & Inside pressure \\
B --- UP 1970--73
  & Falling & Rising fast & Falling fast & Double compression \\
C --- Shock therapy post-1973
  & Slowly rising & Falling & Rising & Denominator destruction \\
\bottomrule
\end{tabular}
\caption{Leverage ratio dynamics under three historical configurations.
Scenario~A: the constraint tightens from the liability side alone.
Scenario~B: reserve drain and monetary expansion compound.
Scenario~C: leverage recovers through denominator compression, not
reserve accumulation.}
\label{app:tab:scenarios}
\end{table}

Scenario~A demonstrates the inside-pressure channel in isolation: even
with a stable or positive external position, fiscal monetisation expands
$M_t$ and the leverage ratio falls. Scenario~B is the double compression
that characterises the Unidad Popular period: high $g_{M,t}$ from
wage-fiscal expansion compounds with reserve drainage through the current
account, so that the constraint became acute through the multiplier before
the full reserve crisis materialised. Scenario~C shows the reverse: the
post-1973 austerity programme restored the leverage ratio not through
reserve accumulation but through denominator compression---the contraction
of domestic monetary liabilities.

The multiplier's asymmetric structure across leverage states is summarised
in Table~\ref{app:tab:multiplier}.

\begin{table}[h]
\centering
\small
\begin{tabular}{lll}
\toprule
LevR level & Multiplier
  $\widetilde{\text{LevR}}_t \cdot g_{M,t}$ & Consequence \\
\midrule
High   & Large    & Large absolute drain---from safe position \\
Medium & Moderate & Moving toward constraint \\
Low    & Small    & Small absolute drain---potentially decisive
                    near threshold \\
\bottomrule
\end{tabular}
\caption{Multiplier structure across leverage states.}
\label{app:tab:multiplier}
\end{table}

The state-dependent local projections specification in
Section~\ref{sec:C2_bop_dynamics} exploits precisely this asymmetry:
identical terms-of-trade shocks are expected to produce different domestic
impulse responses depending on the leverage state at the moment the shock
arrives.
